\documentclass[12pt,a4paper]{article}
\usepackage[margin=2.5cm]{geometry}
\usepackage{amsmath,amssymb,amsthm}
\usepackage{hyperref}
\usepackage{enumitem}

\newtheorem{theorem}{Theorem}[section]
\newtheorem{lemma}[theorem]{Lemma}
\newtheorem{conjecture}[theorem]{Conjecture}
\newtheorem{definition}[theorem]{Definition}
\newtheorem{proposition}[theorem]{Proposition}
\newcommand{\R}{\mathbb{R}}
\newcommand{\C}{\mathbb{C}}
\newcommand{\Z}{\mathbb{Z}}
\newcommand{\N}{\mathbb{N}}
\newcommand{\dist}{\mathcal{D}'}
\newcommand{\tors}{\operatorname{tors}}
\newcommand{\supp}{\operatorname{supp}}

\title{\bf A Greedy Harmonic Substitution in the Geometric Framework\\ for the Riemann Hypothesis}
\author{}
\date{}

\begin{document}
\maketitle

\begin{abstract}
We propose a research programme that modifies the geometric approach to the Riemann Hypothesis (RH) by replacing the smooth trigonometric functions in the Riemann helix with their \emph{greedy harmonic sine reconstruction} counterparts.  The resulting discrete helix has a torsion that is exactly a sum of Dirac delta distributions supported at scaled thresholds of the greedy algorithm.  By matching these thresholds to the zero-counting function of the Riemann zeta function, we aim to prove the Geometric Explicit Formula (GEF) — that the torsion equals the unconditional prime distribution plus a smooth background — using purely combinatorial and asymptotic analysis of the greedy sieve.  This substitution transforms the GEF from an analytic conjecture into a problem about the correlation kernel of the greedy harmonic decomposition, potentially opening a rigorous path to RH.
\end{abstract}

\section{Introduction}

The geometric programme for the Riemann Hypothesis, developed in earlier work, constructs two families of space curves: the twin conical helices (yielding an unconditional prime distribution \(\tau_{\mathrm{prime}}\)) and the \emph{Riemann helix} \(C_1(t)=(\cos\phi(t),\sin\phi(t),t)\).  The phase function \(\phi\) is chosen so that \(\phi(\gamma_n)=2\pi n\) for the imaginary parts \(\gamma_n\) of the zeros of \(\zeta(s)\).  The central conjecture, the \emph{Geometric Explicit Formula} (GEF), states that the torsion of \(C_1\) coincides, in the distributional sense, with \(\tau_{\mathrm{prime}}\) plus a smooth archimedean term.  It is equivalent to RH.

In the present programme we replace the standard sine and cosine by their deterministic \emph{greedy harmonic sine reconstruction} — a piecewise-constant approximation of \(\sin\theta\) on \([0,\pi]\) built from harmonic angles \(\alpha_n=\pi/(n+2)\).  This substitution turns the smooth Riemann helix into a polygonal curve whose torsion singularities are located explicitly at the rescaled thresholds \(\theta_n^*\).  The distribution of these thresholds is governed by a combinatorial recurrence.  We propose to prove that, when the phase \(\phi\) is derived from the zero-counting function via the greedy harmonic sieve, the torsion measure coincides exactly with \(\tau_{\mathrm{prime}} + \tau_{\mathrm{arch}}\), thereby establishing the GEF and, conditionally, RH.

\section{The Geometric Framework (Recap)}

We briefly recall the main objects.  Let \(s=\sigma+it\) with \(\sigma>1\).  The twin conical helices
\[
\mathcal{C}_\pm(x) = \bigl( x^{-\sigma}\cos(t\log x),\; \pm x^{-\sigma}\sin(t\log x),\; x \bigr),\qquad x\ge1,
\]
and their M\"obius-filtered superposition (with the M\"obius function \(\mu\))
\[
\mathcal{C}_\mu(x)=\sum_{n=1}^\infty\frac{\mu(n)}{n}\bigl[ \mathcal{C}_+(nx)+\mathcal{C}_-(nx) \bigr]
\]
produce, after a limit, the unconditional geometric prime distribution
\[
\tau_{\mathrm{prime}}(t)=\sum_{p,m\ge1}\frac{\log p}{p^{m/2}}\,\delta(t-m\log p). \tag{1}
\]

Separately, let \(\rho_n=\beta_n+i\gamma_n\) (\(\gamma_n>0\)) be the non-trivial zeros of \(\zeta(s)\).  Choose a smooth strictly increasing \(\phi:[\gamma_1,\infty)\to\R\) with \(\phi(\gamma_n)=2\pi n\).  The unit Riemann helix
\[
C_1(t)=(\cos\phi(t),\sin\phi(t),t)
\]
has torsion
\[
\tors_1(t)=\frac{\phi'(t)[\phi'(t)^4-\phi'(t)\phi'''(t)+3\phi''(t)^2]}{(\phi'(t)^2+1)(\phi'(t)^4+\phi''(t)^2+\phi'(t)^6)}. \tag{2}
\]

\begin{conjecture}[Geometric Explicit Formula]\label{conj:GEF}
There exists a phase \(\phi\) such that, in \(\dist(\R)\),
\[
\tors_1 = \tau_{\mathrm{prime}} + \tau_{\mathrm{arch}},
\]
with \(\tau_{\mathrm{arch}}\) smooth and explicit.
\end{conjecture}

\section{Greedy Harmonic Sine Reconstruction}

The greedy harmonic sine approximation of \(\sin\theta\) on \([0,\pi]\) is the function
\[
S_N(\theta)=\sum_{n=0}^{N-1}\alpha_n\,\delta_n(\theta),\qquad
\alpha_n=\frac{\pi}{n+2},
\]
where each \(\delta_n(\theta)=\mathbf{1}_{[\theta_n^*,\pi]}(\theta)\) is a step function.  The thresholds \(\theta_n^*\) are determined by the fixed-point equation
\[
\theta_n^* = \alpha_n + \frac{1}{\pi}\sum_{m=0}^{n-1}\alpha_m\theta_m^*, \qquad \theta_{-1}^*:=0. \tag{3}
\]
Only those \(n\) with \(\theta_n^*<\pi\) are \emph{active}; their set is denoted \(\mathcal{A}\).  For \(\theta\in[0,\pi]\) the convergence
\[
\lim_{N\to\infty}S_N(\theta)=\sin\theta
\]
is uniform.  The error is \(O(1/N)\).

The corresponding approximation for \(\cos\theta\) is obtained by the shift \(\theta\mapsto\pi/2-\theta\), which acts on the thresholds in a predictable way.

\section{The Discrete Riemann Helix}

We now define a \emph{discrete} helix by replacing \(\cos\phi(t)\) and \(\sin\phi(t)\) with their greedy harmonic counterparts.  To do so, we first construct a phase function \(\phi_G(t)\) from the zero-counting function \(N(t)\) using the greedy sieve.

\subsection{Greedy phase function}
Recall that \(N(t)\sim \frac{t}{2\pi}\log\frac{t}{2\pi}\).  Define the smooth monotone function
\[
u(t)=\pi\,N(t).
\]
For large \(t\), \(u(t)\approx \frac{t}{2}\log\frac{t}{2\pi}\).  The inverse \(t(u)\) is approximately \(2u/\log u\).  The zeros are located at \(u_n=u(\gamma_n)=2\pi n\).  The greedy harmonic sine reconstruction \(S_N(\theta)\) approximates \(\sin\theta\).  Replacing \(\theta\) by a scaled version of \(u\) yields a candidate for the phase.  More precisely, define
\[
\phi_G(t) = S_{N(t)}\!\Bigl( \frac{u(t)}{R} \Bigr),
\]
where \(R\) is a large scaling factor chosen so that the active thresholds \(\theta_n^*\) map to the zero ordinates.  As \(N(t)\to\infty\) with \(t\), the approximation becomes exact.

\subsection{Discrete helix}
The \emph{greedy harmonic Riemann helix} is the piecewise-linear curve
\[
\widetilde{C}_1(t) = \bigl( \cos_G(\phi_G(t)),\; \sin_G(\phi_G(t)),\; t \bigr),
\]
where \(\cos_G\) and \(\sin_G\) are the greedy harmonic approximations of cosine and sine, respectively.  Because \(\sin_G\) and \(\cos_G\) are step functions of the argument, \(\widetilde{C}_1(t)\) is a polygonal curve with vertices at the points where \(\phi_G(t)\) equals a threshold \(\theta_n^*\).  Its torsion is then a sum of Dirac delta distributions supported at those vertices, weighted by the jumps in the unit tangent and binormal vectors.

\section{Torsion of the Discrete Helix}

For a polygonal space curve, the torsion is concentrated at the vertices.  A careful Frenet–Serret analysis for piecewise-linear curves (using limiting smooth interpolations) shows that if the vertex angles are small, the torsion singularity at a vertex is proportional to the dihedral angle between successive osculating planes.  In our case, the vertices occur at times \(t_k\) where \(\phi_G(t_k)=\theta_{a_k}^*\), and the weights are controlled by the correlation kernel
\[
K(n,m)=\int_0^\pi \bigl(\delta_n(\theta)-\tfrac{\theta}{\pi}\bigr)\bigl(\delta_m(\theta)-\tfrac{\theta}{\pi}\bigr)\,d\theta .
\]

We propose to prove that, when the scaling \(R\) is chosen appropriately and \(N\to\infty\), the distributional limit of \(\tors_{\widetilde{C}_1}\) is exactly
\[
\sum_{k} \frac{\alpha_{a_k}}{\Delta t_k}\,\delta(t-t_k) \;+\; \text{smooth terms},
\]
and that under the map \(k\mapsto m\log p\) (the Greedy–Prime Correspondence) this sum becomes \(\tau_{\mathrm{prime}}\).  The key point is that the harmonic weights \(\alpha_n\) are exactly \(\pi/(n+2)\), which after rescaling turn into the von Mangoldt weights \(\log p/p^{m/2}\).  The smooth background \(\tau_{\mathrm{arch}}\) emerges from the continuous interpolation between vertices.

\section{A Rigorous Programme}

The following steps would constitute a proof of the GEF (and hence RH) using this method:

\begin{enumerate}[label=(\arabic*)]
\item \textbf{Explicit torsion of a polygonal helix.}  
Derive a general formula for the torsion singularities of a piecewise-linear curve \(C(t)=(\cos\Phi(t),\sin\Phi(t),t)\) where \(\Phi\) is a non-decreasing step function with jumps at \(t_k\).  Show that the singular part of \(\tors\) is
\[
\sum_k c_k\,\delta(t-t_k),\qquad c_k = \frac{\Delta\Phi'(t_k)}{\text{something}} .
\]

\item \textbf{Greedy phase function from zero density.}  
Construct \(\Phi(t)\) by scaling the greedy harmonic sine reconstruction to match the zero-counting function.  More specifically, set
\[
\Phi'(t) \approx \frac{2\pi}{R} S'_{N}\!\Bigl(\frac{u(t)}{R}\Bigr)\, u'(t),
\]
so that the jumps of \(\Phi'\) occur at the rescaled thresholds \(t_k\) and have heights proportional to \(\alpha_n\).  Show that the resulting \(\Phi\) is strictly increasing and satisfies the zero-locking condition \(\Phi(\gamma_n)=2\pi n\) asymptotically.

\item \textbf{Kernel-weighted sum to prime weights.}  
Using the correlation kernel \(K(n,m)\) and the property \(\sum_n \alpha_n \delta_n(\theta)\to\sin\theta\), prove that the cumulative sum of the step sizes satisfies
\[
\sum_{t_k\le T} \Delta\Phi(t_k) \sim \sum_{m\log p\le T} \frac{\log p}{p^{m/2}},
\]
with an error that can be controlled by the uniform convergence rate of the greedy sine approximation.  This is the central analytic step, replacing the Weil explicit formula by the greedy harmonic combinatorial identity.

\item \textbf{Archimedean remainder.}  
Show that the smooth part of the torsion, obtained from the continuous interpolation of the step function \(\Phi\), equals \(\tau_{\mathrm{arch}}(t)\).  This relies on standard asymptotics of the zero-counting function.

\item \textbf{Limit and distributional equality.}  
Pass to the limit \(N\to\infty\) (width of the greedy approximation) and mollifier width \(\varepsilon\to0\) to obtain \(\tors_{\widetilde{C}_1} = \tau_{\mathrm{prime}} + \tau_{\mathrm{arch}}\) in \(\dist(\R)\).
\end{enumerate}

\section{Advantages of the Greedy Substitution}

\begin{itemize}
\item \textbf{Exact singularities.}  The polygonal helix has explicit Dirac deltas at known positions, removing the need for an uncontrolled distributional limit of smooth functions.
\item \textbf{Combinatorial control.}  The weights are directly expressed via the harmonic angles and the correlation kernel, which are explicit rational numbers.
\item \textbf{No explicit formula inversion.}  The usual Weil explicit formula is replaced by the combinatorial identity that the greedy harmonic approximation converges to sine.  The prime distribution emerges naturally from the harmonic series via the mapping \(k\mapsto\log p\).
\item \textbf{Numerical verifiability.}  The entire construction can be tested numerically for large finite \(N\), providing strong evidence before a full proof is attempted.
\end{itemize}

\section{Open Problems and Conjectures}

The programme rests on two major conjectures that must be proved.

\begin{conjecture}[Greedy-Prime Correspondence for Thresholds]
Under the rescaling that maps the zero-counting function to the greedy harmonic sine reconstruction, the active threshold positions \(\{t_k\}\) coincide exactly with the set \(\{m\log p : p\text{ prime}, m\ge1\}\), and the jump sizes \(\Delta\Phi(t_k)\) equal \(\log p/p^{m/2}\).
\end{conjecture}

\begin{conjecture}[Polygonal Torsion Limit]
The torsion of the polygonal greedy harmonic helix converges in \(\dist(\R)\) to the sum of Dirac deltas with the above coefficients plus the smooth archimedean background.
\end{conjecture}

If both are proved, the GEF follows and thus RH.

\section{Conclusion}
We have outlined a concrete research programme that replaces the smooth trigonometric functions in the geometric framework for RH with their greedy harmonic sine approximations.  This transforms the central conjecture into a problem about the torsion of a piecewise-linear curve, whose singularities are governed by an explicit combinatorial kernel.  The programme reduces RH to proving that the greedy harmonic sieve, when scaled by the zero-counting function, reproduces the prime distribution.  While challenging, this approach offers a new, self-contained, and potentially rigorous path to the Riemann Hypothesis, bypassing the need for heavy analytic number theory.

\end{document}